\chapter[Der beste Führer]{Der beste Führer}
\quellseite{9}
Es ist erst eine kurze Frist, seitdem wir das neue Jahr begonnen haben. Manche von uns haben mit frohen Hoffnungen und frischem Muth, viele mit stillem Bangen den neuen Lebenspfad betreten, der so unbekannt vor uns allen liegt. Jedes hat sich gewiß dabei die Frage vorgelegt: Wird derselbe durch sonnige Auen oder durch dunkle Thäler führen? Wird mein Fuß nicht gleiten und werde ich den rechten Weg nicht verfehlen? Mir ist unter solchen Fragen eine kleine Begebenheit eingefallen.

Während unseres letzten Aufenthalts in der Schweiz wurde ein Ausflug auf eine der herrlichsten Alpenhöhen ausgeführt. „Ohne Führer geht es nicht!“ so wurde uns gleich von vornherein erklärt. Zuerst schien uns ein solcher überflüssig, aber mehr und mehr lernten wir ihn schätzen. Wie sicher ging es unter seiner bewährten Leitung vorwärts. Kam ein Kreuzweg, so blickte alles auf ihn, um seiner Weisung zu folgen. Führte der Pfad über gefährliche Klippen, an steilen Abgründen vorbei, so erfaßte eins nach dem andern des Führers starke Hand. An besonders schmalen Stellen ermahnte er uns, nur immer genau in seine Fußstapfen zu treten, und so überwanden wir alle Schwierigkeiten. Ueberdies durften wir auch unsere Lasten auf seine starken Schultern legen, und wenn da und dort eins müde werden wollte, sprach er uns Muth zu, indem er uns das baldige Ziel in Aussicht stellte. Endlich war die Höhe erreicht, der Ausblick in die wunderbar schöne Alpenwelt lohnte reichlich alle Mühe und jeden sauren Tritt. Wir mußten die Hände falten und an jenes noch viel herrlichere Ziel denken, das uns einst dort oben erwartet, wenn unsere Erdenwallfahrt vollbracht sein wird.

Und, Gott Lob, auch auf dieser großen Wanderschaft will ein treuer Führer uns zur Seite stehen. O daß ein jedes sich diesen erwählen möchte! Kein irdischer Führer kann uns so behüten wie der, von dem wir bekennen müssen: „Er erquicket meine Seele, er führet mich auf rechter Straße um seines Namens willen.“ Dürfen wir nicht alle unsere Sorgen und Nöthe auf ihn werfen? Stärkt er uns nicht die müden Kniee? Ja, seht nur hinein in eure Bibel: jede köstliche Verheißung, die uns dieser treueste aller Führer im Alten und Neuen Testament gibt, wird sich in seiner Nachfolge täglich an uns erfüllen. Und wenn zuletzt der Weg immer steiler wird, wenn uns der Muth sinken und die Kraft verlassen will, dann sollen wir erst recht erfahren, daß sein Wort an uns zur Wahrheit wird: „Ich will dich tragen bis ins Alter.“ Ihr, die ihr von der langen, oft beschwerlichen Pilgerreise müde und, wie es je und je von den Vätern im alten Bund heißt, lebenssatt seid, ihr dürft immer wieder aufs neue rühmen:

\begin{verse}
Der Alten und der Müden\\
Gedenkt du, Herr, so treu,\\
Dein Wort bringt ihnen Frieden,\\
Erquicket sie täglich neu.

Du sprichst: „Ich will dich heben.“\\
Ob höher steigt die Fluth,\\
Soll doch mein Herz nicht beben,\\
Ich bleib in sicherer Hut.

Du sprichst: „Ich will dich tragen.“\\
Wenn mich die Kraft verläßt,\\
Dann will ich nicht verzagen,\\
Du hältst mich stark und fest.

Du willst mich auch erretten\\
Aus jeder Sklaverei,\\
Du brichst des Todes Ketten\\
Und machst mich herrlich frei.

In deinen treuen Armen\\
Befürcht ich keine Noth,\\
Du trägst mich voll Erbarmen\\
Im Leben und im Tod.
\end{verse}
(Christenbote.)
